\chapter{Conclusiones y Trabajo Futuro}
\label{cap:conclusiones}

\lettrine[lraise=-0.1, lines=2, loversize=0.2]{E}{ste} trabajo de investigación ha abordado el problema de la detección pasiva de vehículos aéreos no tripulados (\gls{uav}) operando en la banda no licenciada ISM de 2.4\,GHz. El escenario de trabajo se ha caracterizado por una atenuación extrema y la presencia de interferencia co-canal severa procedente de redes inalámbricas comerciales. El desafío analítico central ha consistido en procesar señales radioeléctricas afectadas por un nivel crítico de ruido de etiqueta (\textit{label noise}). Este fenómeno deriva de la utilización de ventanas de observación amplias (75\,ms) sobre el canal electromagnético, donde los escasos instantes activos de las ráfagas del dron quedan sumergidos en extensos períodos de ruido térmico; una condición intrínseca a la complejidad del conjunto de datos público NoisyUAV v2 evaluado a lo largo del estudio.

\section{Conclusiones}
\label{sec:conclusiones}

Para solventar la incapacidad demostrada por los detectores de energía convencionales en entornos con relaciones señal a ruido negativas, se ha diseñado un detector analítico de preprocesamiento fundamentado en la entropía espectral. 

Este bloque determinista opera sin requerir parámetros entrenables. Aplica una etapa de blanqueo espectral secuenciada con un umbral \gls{cfar} adaptativo para aislar la morfología de las ráfagas \gls{fhss}. Su validación demuestra que es capaz de localizar transmisiones de drones en regímenes de \gls{snr} donde la magnitud de la señal útil queda matemáticamente enmascarada por el suelo de ruido del receptor.

La evolución del motor de clasificación evidenció rápidamente que las arquitecturas superficiales carecían de resiliencia frente al ruido, mientras que el incremento masivo de la profundidad convolucional (\textit{DualStream-Deep}) provocaba un colapso en la generalización. La solución técnica definitiva se ha consolidado en la topología de redes neuronales convolucionales de valores complejos (\gls{cv-cnn}) denominada \textit{DualStream-DynamicSlidingWindow}. Esta arquitectura procesa de forma nativa la matriz temporal I/Q preservando la coherencia de fase. Simultáneamente, la complementa con una estimación de la densidad espectral de potencia (\gls{psd}) gobernada por un bloque de atención dinámica. El sistema descarta además la evaluación estática; en su lugar, implementa una ventana temporal deslizante que recalcula el umbral de ruido localmente de forma ininterrumpida.

La evaluación experimental del modelo definitivo constata un F1-Score del 90.12\,\% y una Especificidad del 97.10\,\% operando bajo el modo de máxima seguridad, es decir, minimizando los falsos positivos. De forma complementaria, al configurar la arquitectura en su modo de alerta temprana, el sistema relaja el umbral de decisión para maximizar el \textit{Recall}. Esta segunda configuración prioriza la detección precoz de la amenaza y garantiza la intercepción de señales muy por debajo de los 0\,dB, asumiendo un incremento controlado en la tasa de falsas alarmas que resulta idóneo para la vigilancia de perímetros abiertos.

La hibridación de los flujos I/Q y \gls{psd} iguala la precisión publicada en la literatura para arquitecturas de referencia basadas en técnicas clásicas de visión artificial \cite{gluge_robust_2024}, logrando estabilizar la detección por encima de los $-2$\,dB. Este rendimiento equivalente se alcanza requiriendo únicamente 344.291 parámetros entrenables, lo que supone un factor de reducción paramétrica de entre 27 y 58 veces respecto al estado del arte. Se elimina así el coste computacional asociado a la extracción de espectrogramas bidimensionales, posicionando el procesamiento unidimensional complejo como una alternativa técnica superior.

A pesar de estos avances empíricos, el diseño asume ciertas limitaciones metodológicas. La validación se ha concentrado exclusivamente en el repositorio NoisyUAV v2, restringiendo la demostración de generalización frente a otras plataformas aéreas no capturadas en dicho volumen. La ausencia de inferencia sobre hardware de radio definida por software (\gls{sdr}) mantiene la validación en el ámbito algorítmico. Por último, la caída asintótica del rendimiento ante señales del Grupo D (\gls{snr} $< -12$\,dB) ratifica la existencia de un límite físico infranqueable en la recuperación de información muy por debajo del umbral termodinámico del receptor.

\section{Líneas de Trabajo Futuro}
\label{sec:trabajo_futuro}

La validación de la inferencia del modelo sobre hardware embebido constituye la progresión tecnológica directa de este proyecto. A la hora de diseñar las directrices de integración, es obligatorio prescindir de arquitecturas de propósito general basadas en microordenadores ARM convencionales (tipo Raspberry Pi o similares) y pasar a su integración nativa sobre hardware que cuenta con módulos de aceleración nativos, como TPUs embebidas. Algunos ejemplos son plataformas como Google Coral o directamente sobre FPGAs, acopladas a receptores \gls{sdr} comerciales. Para llevar este despliegue a producción, será necesario implementar un ciclo de integración automatizado (generalmente denominado flujo MLOps, \textit{Machine Learning Operations}) que aplique estrategias de monitorización, entrenamiento y despliegue automático, así como prácticas de compresión sobre el modelo. Esto incluye de manera prioritaria el uso de técnicas de cuantificación de pesos post-entrenamiento que reducen los tensores desde punto flotante de 32 bits a representaciones numéricas como INT8 o FP16, con el objetivo de evaluar el balance exacto entre latencia de respuesta (milisegundos por ráfaga), eficiencia energética y la degradación marginal del algoritmo clasificador en un escenario táctico real.

La expansión del espectro de monitorización hacia la banda de 5.8\,GHz representa otra evolución obligatoria.

Esta banda frecuencial concentra los protocolos de transmisión de vídeo FPV predominantes en los dispositivos comerciales. Dado que estos enlaces emplean esquemas de modulación analógicos y flujos digitales continuos muy distintos a las ráfagas intermitentes de telemetría, resulta imprescindible evaluar la capacidad de la arquitectura \gls{cv-cnn} para asimilar estas nuevas morfologías electromagnéticas.

En el ámbito puramente algorítmico, la optimización del entrenamiento frente al ruido de etiqueta inherente a las señales radioeléctricas resulta prioritaria. Explorar paradigmas de aprendizaje débilmente supervisado, como \textit{Multiple Instance Learning} (MIL), permitiría prescindir del descarte manual de los datos pertenecientes a niveles críticos de \gls{snr}. Esta aproximación facilitaría que la red extrajera patrones latentes directamente de las muestras más degradadas sin comprometer la estabilidad matemática del descenso de gradiente a lo largo de las épocas de entrenamiento.

Finalmente, el despliegue del sistema en áreas urbanas revela la necesidad de complementar la detección de radiofrecuencia con otros sensores físicos. En estos escenarios, las señales electromagnéticas sufren rebotes continuos contra los edificios (desvanecimiento por trayectos múltiples), generando distorsiones que pueden engañar al clasificador y disparar falsas alarmas. Para mitigar este problema, una futura línea de investigación consistiría en implementar un esquema de fusión tardía. Esta técnica cruzaría la probabilidad de detección calculada por el modelo RF con la telemetría de un radar de onda continua (FMCW) o las lecturas de una red de micrófonos. Al contrastar la firma radioeléctrica con la confirmación física del espacio aéreo, el sistema global obtendría la redundancia necesaria para anular casi por completo los falsos positivos derivados de las anomalías del canal.
